% !TEX root = ../main.tex
% ============================================================================
% Owner: Arteaga
% ============================================================================
\section{Introduction}

Depth Anything and similar models supply depth to systems that have no depth sensor. They are used in driving, robotics, augmented reality, and 3D reconstruction, usually as a drop-in component. Every pixel gets a depth value and no measure of reliability: an estimate that is wrong by thirty percent looks exactly like one that is right. The cost of that silence is highest in the conditions where camera depth is least reliable: night, rain on the lens, spray, and low sun. A confidence map lets a downstream consumer weight the depth, discard it, or fall back to another sensor.

This project trains a model to predict that error without touching the depth model. A small convolutional head reads the frozen model's internal features along with its predicted depth, the temporal spread of depth across neighboring frames, a radar residual, and a radar validity mask. It outputs, per cell of a $132 \times 74$ grid, the probability that the depth there is wrong by more than a threshold; the confidence is one minus that probability. Lidar projected into the image gives the true error at the cells that receive a return, so the head learns from scattered labels and produces a dense map at inference. Training needs lidar. The design goal is that inference does not, while radar, which a production car does carry, is present in both. So far the heads meet that goal only in part. The temporal spread is computed with a lidar fit of the depth scale, so a head that reads it still depends on lidar; a head that leaves it out, such as the one that reads the decoder features alone, does not (Section~\ref{sec:runs}).
